
\documentclass{article}
\usepackage{colortbl}
\usepackage{makecell}
\usepackage{multirow}
\usepackage{supertabular}

\begin{document}

\newcounter{utterance}

\centering \large Interaction Transcript for game `hot\_air\_balloon', experiment `air\_balloon\_survival\_en\_complexity\_easy', episode 1 with Qwen3{-}32B{-}t0.0.
\vspace{24pt}

{ \footnotesize  \setcounter{utterance}{1}
\setlength{\tabcolsep}{0pt}
\begin{supertabular}{c@{$\;$}|p{.15\linewidth}@{}p{.15\linewidth}p{.15\linewidth}p{.15\linewidth}p{.15\linewidth}p{.15\linewidth}}
    \# & $\;$A & \multicolumn{4}{c}{Game Master} & $\;\:$B\\
    \hline

    \theutterance \stepcounter{utterance}  
    & & \multicolumn{4}{p{0.6\linewidth}}{
        \cellcolor[rgb]{0.9,0.9,0.9}{
            \makecell[{{p{\linewidth}}}]{
                \texttt{\tiny{[P1$\langle$GM]}}
                \texttt{You are participating in a collaborative negotiation game.} \\
\\ 
\texttt{Together with another participant, you must agree on a single set of items that will be kept. Each of you has your own view of how much each item matters to you (importance). You do not know how the other participant values the items. Additionally, you are given the effort each item demands.} \\
\texttt{You may only agree on a set if the total effort of the selected items does not exceed a shared limit:} \\
\\ 
\texttt{LIMIT = 8170} \\
\\ 
\texttt{Here are the individual item effort values:} \\
\\ 
\texttt{Item effort = \{"B21": 594, "B00": 551, "A08": 634, "A27": 204, "B94": 557, "B53": 440, "A23": 734, "C65": 893, "A89": 246, "B67": 592, "B45": 143, "C34": 568, "C12": 472, "A05": 401, "B60": 729, "C92": 201, "B69": 85, "A84": 641, "C06": 958, "B12": 80, "A16": 157, "B85": 806, "C35": 684, "B10": 59, "C15": 31, "B09": 765, "B16": 416, "C81": 392, "A58": 428, "C24": 699, "A30": 141, "A66": 606, "B75": 612, "B46": 133, "C70": 689\}} \\
\\ 
\texttt{Here is your personal view on the importance of each item:} \\
\\ 
\texttt{Item importance values = \{"B21": 138, "B00": 583, "A08": 868, "A27": 822, "B94": 783, "B53": 65, "A23": 262, "C65": 121, "A89": 508, "B67": 780, "B45": 461, "C34": 484, "C12": 668, "A05": 389, "B60": 808, "C92": 215, "B69": 97, "A84": 500, "C06": 30, "B12": 915, "A16": 856, "B85": 400, "C35": 444, "B10": 623, "C15": 781, "B09": 786, "B16": 3, "C81": 713, "A58": 457, "C24": 273, "A30": 739, "A66": 822, "B75": 235, "B46": 606, "C70": 105\}} \\
\\ 
\texttt{Goal:} \\
\\ 
\texttt{Your goal is to negotiate a shared set of items that benefits you as much as possible (i.e., maximizes total importance to YOU), while staying within the LIMIT. You are not required to make a PROPOSAL in every message {-} you can simply negotiate as well. All tactics are allowed!} \\
\\ 
\texttt{Interaction Protocol:} \\
\\ 
\texttt{You may only use the following structured formats in a message:} \\
\\ 
\texttt{PROPOSAL: \{'A', 'B', 'C', …\}} \\
\texttt{Propose keeping exactly those items.} \\
\texttt{REFUSE: \{'A', 'B', 'C', …\}} \\
\texttt{Explicitly reject opponent's proposal.} \\
\texttt{ARGUMENT: \{'...'\}} \\
\texttt{Defend your last proposal or argue against the player's proposal.} \\
\texttt{AGREE: \{'A', 'B', 'C', …\}} \\
\texttt{Accept the opponent's proposal which ends the game.} \\
\texttt{STRATEGIC REASONING: \{'...'\}} \\
\texttt{	Describe your strategic reasoning or anticipation explaining your choice of action. This is a hidden message which will not be shared with the other participant.} \\
\\ 
\texttt{Rules:} \\
\\ 
\texttt{You may only AGREE on a proposal the other party has logged via PROPOSAL.} \\
\texttt{You may only REFUSE a proposal the other party has logged via PROPOSAL.} \\
\texttt{Total effort of any PROPOSAL or AGREE set must be ≤ LIMIT.} \\
\texttt{Do NOT reveal your hidden importance scores.} \\
\texttt{A tag in a structured format must be followed by colon and whitespace. The argument must be a python set containing 0 or more strings.} \\
\texttt{So, it must be of the form TAG: \{...\}} \\
\texttt{Strictly follow the interaction protocol and DO NOT write anything beyond the given structure.} \\
\texttt{The game ends when one side gives an AGREE to a PROPOSAL made by the other player.} \\
\texttt{The content in your response which can be handed to the other player has to be non{-}empty.} \\
\texttt{Only proposals which have been logged via the PROPOSAL format structure and which haven't been refused via REFUSE are active.} \\
\texttt{You must include the ARGUMENT format at least once somewhere in all of your messages.} \\
\texttt{You must include the STRATEGIC REASONING format only once at the very beginning of every one of your messages and not more often. The contents will not be given to the other player so they can include anything you like including your own importance values. Here you should reason step by step to come up with you next move.} \\
            }
        }
    }
    & \\ \\

    \theutterance \stepcounter{utterance}  
    & \multicolumn{4}{p{0.6\linewidth}}{
        \cellcolor[rgb]{0.9,0.9,0.9}{
            \makecell[{{p{\linewidth}}}]{
                \texttt{\tiny{[P1$\rangle$GM]}}
                \texttt{<think>} \\
\texttt{Okay, let's see. I need to figure out the best set of items to propose that maximizes my total importance without exceeding the effort limit of 8170. First, I should calculate the total effort of all items to see if it's over the limit. Let me add them up.} \\
\\ 
\texttt{Looking at the item efforts: B21 is 594, B00 is 551, A08 is 634, A27 is 204, B94 is 557, B53 is 440, A23 is 734, C65 is 893, A89 is 246, B67 is 592, B45 is 143, C34 is 568, C12 is 472, A05 is 401, B60 is 729, C92 is 201, B69 is 85, A84 is 641, C06 is 958, B12 is 80, A16 is 157, B85 is 806, C35 is 684, B10 is 59, C15 is 31, B09 is 765, B16 is 416, C81 is 392, A58 is 428, C24 is 699, A30 is 141, A66 is 606, B75 is 612, B46 is 133, C70 is 689.} \\
\\ 
\texttt{Adding all these up... Let me do that step by step. Maybe group them to make it easier. Let's see:} \\
\\ 
\texttt{Start with B21 (594) + B00 (551) = 1145. Then A08 (634) adds to 1779. A27 (204) brings it to 1983. B94 (557) is 2540. B53 (440) is 2980. A23 (} \\
            }
        }
    }
    & & \\ \\

    \theutterance \stepcounter{utterance}  
    & & \multicolumn{4}{p{0.6\linewidth}}{
        \cellcolor[rgb]{0.9,0.9,0.9}{
            \makecell[{{p{\linewidth}}}]{
                \texttt{\tiny{[P1$\langle$GM]}}
                \texttt{Your response did not start with the proper strategic reasoning tag at the very beginning of your response. The very first structured format must be of the form STRATEGIC REASONING: \{...\}. Try again.} \\
            }
        }
    }
    & \\ \\

    \theutterance \stepcounter{utterance}  
    & \multicolumn{4}{p{0.6\linewidth}}{
        \cellcolor[rgb]{0.9,0.9,0.9}{
            \makecell[{{p{\linewidth}}}]{
                \texttt{\tiny{[P1$\rangle$GM]}}
                \texttt{<think>} \\
\texttt{Okay, I need to make sure I start with the STRATEGIC REASONING tag. Let me think through this again.} \\
\\ 
\texttt{First, I should calculate the total effort of all items to see if the sum is over the limit. If it is, I can't take all. Let me add up all the effort values. } \\
\\ 
\texttt{Looking at the list: B21 (594) + B00 (551) = 1145. Then A08 (634) adds to 1779. A27 (204) makes 1983. B94 (557) is 2540. B53 (440) is 2980. A23 (734) brings it to 3714. C65 (893) is 4607. A89 (246) is 4853. B67 (592) is 5445. B45 (143) is 5588. C34 (568) is 6156. C12 (472) is 6628. A05 (401) is 7029. B60 (729) is 7758. C92 (201) is 7959. B69 (85) is 8044. A84 (641) would push it over 8685, which is way over the limit of 8170. So the total effort of all items is way over. Therefore, I need to select a subset that's under 8170.} \\
\\ 
\texttt{Next, I need to maximize my own importance. So I should prioritize items with high importance per unit effort. Let me calculate the importance per effort ratio for each item. } \\
\\ 
\texttt{For example, B12 has an importance of 915 and effort 80, so 915/80 = 11.44. C15 is 781/31 ≈25.19. B10 is 623/59 ≈10.56. A16 is 8} \\
            }
        }
    }
    & & \\ \\

    \theutterance \stepcounter{utterance}  
    & & \multicolumn{4}{p{0.6\linewidth}}{
        \cellcolor[rgb]{0.9,0.9,0.9}{
            \makecell[{{p{\linewidth}}}]{
                \texttt{\tiny{[P1$\langle$GM]}}
                \texttt{Your response did not start with the proper strategic reasoning tag at the very beginning of your response. The very first structured format must be of the form STRATEGIC REASONING: \{...\}. Try again.} \\
            }
        }
    }
    & \\ \\

    \theutterance \stepcounter{utterance}  
    & \multicolumn{4}{p{0.6\linewidth}}{
        \cellcolor[rgb]{0.9,0.9,0.9}{
            \makecell[{{p{\linewidth}}}]{
                \texttt{\tiny{[P1$\rangle$GM]}}
                \texttt{<think>} \\
\texttt{Okay, I need to make sure I start with the STRATEGIC REASONING tag. Let me think through this again.} \\
\\ 
\texttt{First, I should calculate the total effort of all items to see if the sum is over the limit. If it is, I can't take all. Let me add up all the effort values. } \\
\\ 
\texttt{Looking at the list: B21 (594) + B00 (551) = 1145. Then A08 (634) adds to 1779. A27 (204) makes 1983. B94 (557) is 2540. B53 (440) is 2980. A23 (734) brings it to 3714. C65 (893) is 4607. A89 (246) is 4853. B67 (592) is 5445. B45 (143) is 5588. C34 (568) is 6156. C12 (472) is 6628. A05 (401) is 7029. B60 (729) is 7758. C92 (201) is 7959. B69 (85) is 8044. A84 (641) would push it over 8685, which is way over the limit of 8170. So the total effort of all items is way over. Therefore, I need to select a subset that's under 8170.} \\
\\ 
\texttt{Next, I need to maximize my own importance. So I should prioritize items with high importance per unit effort. Let me calculate the importance per effort ratio for each item. } \\
\\ 
\texttt{For example, B12 has an importance of 915 and effort 80, so 915/80 = 11.44. C15 is 781/31 ≈25.19. B10 is 623/59 ≈10.56. A16 is 8} \\
            }
        }
    }
    & & \\ \\

    \theutterance \stepcounter{utterance}  
    & & & \multicolumn{2}{p{0.3\linewidth}}{
        \cellcolor[rgb]{0.9,0.9,0.9}{
            \makecell[{{p{\linewidth}}}]{
                \texttt{\tiny{[GM$|$GM]}}
                \texttt{abort} \\
            }
        }
    }
    & & \\ \\

    \theutterance \stepcounter{utterance}  
    & & & \multicolumn{2}{p{0.3\linewidth}}{
        \cellcolor[rgb]{0.9,0.9,0.9}{
            \makecell[{{p{\linewidth}}}]{
                \texttt{\tiny{[GM$|$GM]}}
                \texttt{end game} \\
            }
        }
    }
    & & \\ \\

\end{supertabular}
}

\end{document}
